\chapter{Deterministic Cybernetic Systems}
\label{chap:deterministic-cybernetics}

\section{Orientation: Why Cybernetics Belongs Here}
\label{sec:cybernetic-orientation}

The preceding chapters developed a categorical theory of stateful open systems.  Interfaces were
represented by internal categories, stateful processes by Mealy morphisms, comparisons by Mealy
simulations, behaviour by residual semantics, specifications by subobjects and dynamic modalities,
parallelism by product structure, true concurrency by cubical execution, and feedback by hidden-interface
trace.

This chapter gives a concrete application of that machinery to deterministic cybernetics.  The
guiding idea is that a deterministic cybernetic system consists of a plant, an adaptive controller,
an optimizing decision mechanism, and a learning update mechanism connected through feedback.

The naive systems diagram is
\[
\text{plant}
\;\longleftrightarrow\;
\text{controller}
\;\longleftrightarrow\;
\text{policy}
\;\longleftrightarrow\;
\text{optimizer}
\;\longleftrightarrow\;
\text{updater}.
\]
This diagram is useful as informal intuition, but it is too flat as mathematics.  The formal
decomposition used in this chapter is
\[
\boxed{
\text{plant}
\quad\text{feedback-composed with}\quad
\text{adaptive controller}.
}
\]
The adaptive controller then has internal state containing:
\[
\boxed{
\text{controller state}
+
\text{policy/model state}
+
\text{optimizer state}
+
\text{updater state}.
}
\]

The feedback construction supplies the external wiring.  The adaptive controller supplies the
internal decision and learning dynamics.  The resulting closed system is a deterministic process on
the coupled state object.

The word ``deterministic'' is important.  This chapter does not develop stochastic control or
Markov decision processes.  Rather, it identifies the deterministic skeleton underlying such systems:
state changes are functions, and the corresponding Markov kernels are Dirac kernels.  Thus ordinary
Mealy machines are read as deterministic controlled Markov processes with output.

There is also an important terminological convention.  We shall not describe the closed-loop system
as globally clocked.  A global clocking metaphor risks suggesting that independent events have been
reduced to interleaving schedules.  That would be the wrong reading in the presence of the true
concurrency theory developed earlier.  Instead, we speak of a causal Mealy step, a state-mediated
feedback step, or guarded adaptive feedback.  These phrases indicate only that some data depend on
previously produced data.  For example, an updater cannot learn from a reward before the reward has
been produced.  This is causal dependence, not interleaving semantics.

The central thesis of the chapter is therefore:
\[
\boxed{
\text{A deterministic cybernetic system is a state-mediated feedback composite of a plant and
an adaptive controller.}
}
\]
Equivalently, after closing the feedback interface, it is a deterministic Markov process on the
coupled state object.

\section{Deterministic Mealy Processes as Dirac Markov Processes}
\label{sec:mealy-dirac-markov}

We begin with the elementary bridge from Mealy machines to deterministic Markov processes.

\begin{definition}[Deterministic Mealy process]
\label{def:deterministic-mealy-process}
A deterministic Mealy process from an input set \(I\) to an output set \(O\) consists of a state
set \(S\), a transition map
\[
\tau:S\times I\longrightarrow S,
\]
and an output map
\[
\omega:S\times I\longrightarrow O.
\]
Given a current state \(s_t\in S\) and an input \(i_t\in I\), execution is given by
\[
s_{t+1}=\tau(s_t,i_t),
\qquad
o_t=\omega(s_t,i_t).
\]
\end{definition}

For each fixed input \(i\in I\), the transition map
\[
\tau_i:S\longrightarrow S,
\qquad
\tau_i(s)=\tau(s,i),
\]
determines a deterministic Markov kernel
\[
K_i:S\longrightarrow \mathcal D(S)
\]
by
\[
K_i(s)=\delta_{\tau(s,i)},
\]
where \(\mathcal D(S)\) denotes the set of probability distributions on \(S\), and
\(\delta_x\) is the Dirac distribution concentrated at \(x\).

If output is retained as part of the next observable event, we obtain the Dirac kernel
\[
\overline K_i:S\longrightarrow \mathcal D(S\times O)
\]
defined by
\[
\overline K_i(s)=\delta_{(\tau(s,i),\omega(s,i))}.
\]

Thus a deterministic Mealy process is exactly a controlled Markov process whose transition
probabilities take only the values \(0\) and \(1\).

\begin{proposition}[Mealy processes and Dirac Markov processes]
\label{prop:mealy-dirac-markov}
Every deterministic Mealy process determines a Dirac-valued controlled Markov process with output.
Conversely, every Dirac-valued controlled Markov process with output determines a deterministic
Mealy process.
\end{proposition}

\begin{proof}
Given a deterministic Mealy process \((S,\tau,\omega)\), define
\[
\overline K_i(s)=\delta_{(\tau(s,i),\omega(s,i))}.
\]
This is a controlled Markov kernel whose value at each state is a Dirac distribution.

Conversely, suppose that for each input \(i\in I\) we have a Markov kernel
\[
\overline K_i:S\longrightarrow \mathcal D(S\times O)
\]
such that \(\overline K_i(s)\) is always a Dirac distribution.  Then for each \(s\in S\) and
\(i\in I\), there is a unique pair \((s',o)\in S\times O\) such that
\[
\overline K_i(s)=\delta_{(s',o)}.
\]
Define
\[
\tau(s,i)=s',
\qquad
\omega(s,i)=o.
\]
This gives a deterministic Mealy process.  The two constructions are inverse to one another.
\end{proof}

The categorical point is that the Mealy process already contains the Markov state.  The Markov
property says that the future depends on the past only through the current state.  In a deterministic
Mealy process, this dependence is functional rather than probabilistic.

In the one-object specialization of the general theory, a Mealy morphism is a state object with a
right action of the input monoid and an output cocycle valued in the output monoid.  The Dirac
kernel description is the corresponding finite-set probabilistic shadow.

\section{Plants, Observations, Rewards, and Action Interfaces}
\label{sec:plants-observations-rewards}

A cybernetic system contains a plant: the system being regulated.

\begin{definition}[Deterministic plant]
\label{def:deterministic-plant}
A deterministic plant consists of:
\[
X=\text{plant state object},
\]
\[
A=\text{action interface},
\]
\[
O=\text{observation interface},
\]
\[
R=\text{reward, cost, or performance interface},
\]
together with maps
\[
p:X\times A\longrightarrow X,
\]
\[
h:X\longrightarrow O,
\]
and
\[
\rho:X\times A\times X\longrightarrow R.
\]
Given a state \(x\in X\) and an action \(a\in A\), the next plant state is
\[
x'=p(x,a),
\]
the resulting observation is
\[
o'=h(x'),
\]
and the resulting performance signal is
\[
r=\rho(x,a,x').
\]
\end{definition}

The plant may be packaged as a Mealy process
\[
P:A\rightsquigarrow O\times R
\]
with internal state \(X\).  Its transition-output structure is
\[
P(x,a)=
\bigl(
p(x,a),\,h(p(x,a)),\,\rho(x,a,p(x,a))
\bigr).
\]

The convention in this chapter is that the plant emits the post-action observation and performance
signal.  Thus an action produces a plant transition, and the resulting state is then observed:
\[
a_t
\leadsto
x_{t+1}
\leadsto
(o_{t+1},r_t).
\]
The controller may also use a current observation \(o_t=h(x_t)\) to choose the action.  The
post-action observation \(o_{t+1}\) is then used by the adaptive updater.

This distinction prevents the following two maps from being conflated:
\[
h(x_t)=o_t,
\qquad
h(x_{t+1})=o_{t+1}.
\]
The first is used for present decision-making.  The second is used for learning from the consequence
of the decision.

\section{Adaptive Controllers as Stateful Mealy Systems}
\label{sec:adaptive-controllers}

A controller in a deterministic cybernetic system is not merely a function
\[
O\longrightarrow A.
\]
That is only a stateless policy readout.  A learning controller is stateful.

\begin{definition}[Adaptive controller]
\label{def:adaptive-controller}
An adaptive controller from observation/performance interface \(O\times R\) to action/value
interface \(A\times V\) is a deterministic Mealy process
\[
K:O\times R\rightsquigarrow A\times V
\]
whose internal state object is decomposed as
\[
K_0=C\times \Theta\times Z\times U.
\]
Here:
\[
C=\text{controller mode, guard, or safety state},
\]
\[
\Theta=\text{policy, model, parameter, or critic state},
\]
\[
Z=\text{optimizer state},
\]
\[
U=\text{updater or memory state}.
\]
\end{definition}

The controller receives an observation/performance signal and emits an action together with a value
or certificate:
\[
(o,r)
\longmapsto
(a,v).
\]
The selected action \(a\) is sent to the plant.  The value \(v\) is an exposed performance or
optimization certificate.

The policy itself should be separated from the policy state.  For a fixed parameter or model state
\(\theta\in\Theta\), there may be a stateless policy readout
\[
\pi_\theta:O\longrightarrow A.
\]
But the adaptive policy system is stateful because \(\theta\) changes:
\[
\theta\longmapsto\theta'.
\]
Thus the accurate statement is not merely that the policy is stateful.  Rather, the adaptive
controller has policy, model, or critic state.

\section{Policy States, Evaluators, Optimizers, and Updaters}
\label{sec:policy-evaluator-optimizer-updater}

The adaptive controller internally separates four functions:
\[
\text{guards},
\qquad
\text{evaluation},
\qquad
\text{optimization},
\qquad
\text{update}.
\]
This separation is the main structural point of the chapter.

\subsection{Guards and Admissible Actions}
\label{subsec:guards-admissible-actions}

The controller state may restrict the actions currently allowed.  We represent this by a guard
relation
\[
G\hookrightarrow C\times O\times A.
\]
In finite-set notation, this is equivalently a map
\[
G:C\times O\longrightarrow \mathcal P(A),
\]
where \(G(c,o)\subseteq A\) is the set of actions admissible in controller state \(c\) under
observation \(o\).

For example:
\[
\text{a thermostat may forbid rapid switching},
\]
\[
\text{an inventory controller may forbid repeated large orders},
\]
\[
\text{a cruise controller may forbid acceleration in an unsafe following state},
\]
\[
\text{a navigation controller may forbid moves known to be blocked}.
\]

\subsection{Evaluators}
\label{subsec:evaluators}

Let \(W\) be an ordered value object.  An evaluator is a map
\[
Q:\Theta\times C\times O\times A\longrightarrow W.
\]
For fixed \((\theta,c,o)\), the function
\[
a\longmapsto Q(\theta,c,o,a)
\]
scores each candidate action.

The evaluator is the location of value functions, cost functions, predictive models, safety
penalties, learned estimates, and critic values.  It does not itself choose an action.  It only
assigns values.

\subsection{Optimizers}
\label{subsec:optimizers}

The optimizer performs the current decision.  It has state \(Z\), reads the current policy/model
state, controller state, and observation, and returns a selected action, an optimal value, and an
updated optimizer state:
\[
\Omega:Z\times \Theta\times C\times O\longrightarrow A\times W\times Z.
\]
Write
\[
\Omega(z,\theta,c,o)=(a^*,v^*,z^+).
\]
The intended equations are
\[
a^*\in \operatorname{argmax}_{a\in G(c,o)} Q(\theta,c,o,a),
\]
and
\[
v^*=\max_{a\in G(c,o)} Q(\theta,c,o,a).
\]
Thus the optimizer carries both an \(\operatorname{argmax}\) component and a \(\max\) component.
The selected action \(a^*\) controls the plant.  The value \(v^*\) is the value certificate.

\subsection{The Optimizing Subobject}
\label{subsec:optimizing-subobject}

In a more categorical presentation, the optimizer may be represented by an optimizing subobject.
Define
\[
\operatorname{Opt}_Q
\hookrightarrow
Z\times\Theta\times C\times O\times A\times W\times Z.
\]
An element
\[
(z,\theta,c,o,a,v,z^+)
\]
belongs to \(\operatorname{Opt}_Q\) when:
\[
a\in G(c,o),
\]
\[
Q(\theta,c,o,a)=v,
\]
and for every admissible action \(b\in G(c,o)\),
\[
Q(\theta,c,o,b)\leq v.
\]

If the optimizing subobject is represented by a function, then one has a deterministic optimizer
\[
\Omega:Z\times\Theta\times C\times O\longrightarrow A\times W\times Z.
\]
If not, the optimizer is naturally span-valued.  In that case, multiple actions may realize the
same maximum.  This chapter mostly uses deterministic representatives, but the span-valued form is
the natural generalization.

\subsection{Updaters}
\label{subsec:updaters}

The updater performs learning.  It changes the future controller and policy/model state after the
plant consequence is observed.

An updater is a map
\[
L:
C\times\Theta\times Z\times U
\times O\times A\times W\times R\times O
\longrightarrow
C\times\Theta\times Z\times U.
\]
The input variables are:
\[
(c,\theta,z,u,o,a,v,r,o').
\]
Here \(o\) is the observation used for the present decision, \(a\) is the selected action, \(v\)
is the optimizer's value certificate, \(r\) is the resulting performance signal, and \(o'\) is the
post-action observation.

The updater may implement:
\[
\text{table update},
\qquad
\text{model update},
\qquad
\text{safety-state update},
\]
\[
\text{controller-mode update},
\qquad
\text{optimizer-cache update},
\qquad
\text{memory update}.
\]

The central distinction is:
\[
\boxed{
\text{optimizer = within-step decision},
}
\]
whereas
\[
\boxed{
\text{updater = across-step adaptation}.
}
\]
The optimizer chooses using the current policy/model state.  The updater changes the policy/model
state for future choices.

\section{The Causal Closed-Loop Mealy Transition}
\label{sec:causal-closed-loop-transition}

We now define the closed-loop deterministic system obtained from a plant and an adaptive controller.

Let the plant be
\[
(X,A,O,R,p,h,\rho),
\]
and let the adaptive controller have state object
\[
C\times\Theta\times Z\times U.
\]
The coupled closed-loop state object is
\[
S=X\times C\times\Theta\times Z\times U.
\]

Given
\[
s=(x,c,\theta,z,u)\in S,
\]
define the current observation
\[
o=h(x).
\]
The optimizer produces
\[
(a^*,v^*,z^+)=\Omega(z,\theta,c,o).
\]
The plant then responds:
\[
x'=p(x,a^*).
\]
The resulting observation is
\[
o'=h(x').
\]
The performance signal is
\[
r=\rho(x,a^*,x').
\]
Finally, the updater produces the next adaptive-controller state:
\[
(c',\theta',z',u')
=
L(c,\theta,z^+,u,o,a^*,v^*,r,o').
\]

\begin{definition}[Closed-loop transition]
\label{def:closed-loop-transition}
The closed-loop transition of the deterministic adaptive cybernetic system is the function
\[
F:X\times C\times\Theta\times Z\times U
\longrightarrow
X\times C\times\Theta\times Z\times U
\]
defined by
\[
F(x,c,\theta,z,u)=(x',c',\theta',z',u'),
\]
where \(x'\), \(c'\), \(\theta'\), \(z'\), and \(u'\) are computed by the preceding equations.
\end{definition}

This is a causal Mealy transition.  The word ``causal'' only means that the update reads the
consequence of the action.  It does not impose a global interleaving semantics on independent
components.

\begin{theorem}[Closed-loop deterministic Markov process]
\label{thm:closed-loop-deterministic-markov}
Every deterministic adaptive cybernetic system determines a deterministic Markov process on the
closed-loop state object
\[
S=X\times C\times\Theta\times Z\times U.
\]
\end{theorem}

\begin{proof}
The displayed construction defines a function
\[
F:S\longrightarrow S.
\]
The corresponding Markov kernel is the Dirac kernel
\[
K(s)=\delta_{F(s)}.
\]
Thus the closed-loop system is a deterministic Markov process.
\end{proof}

\begin{corollary}[Eventually periodic finite behaviour]
\label{cor:eventually-periodic-cybernetic}
If \(X\), \(C\), \(\Theta\), \(Z\), and \(U\) are finite, then every initial closed-loop state has
eventually periodic state behaviour under \(F\).
\end{corollary}

\begin{proof}
The state object
\[
S=X\times C\times\Theta\times Z\times U
\]
is finite.  Therefore the sequence
\[
s,\ F(s),\ F^2(s),\ldots
\]
eventually repeats.
\end{proof}

\section{Feedback Closure and Signed Interfaces}
\label{sec:feedback-signed-interfaces-cybernetics}

The preceding section defined the closed-loop execution map.  We now relate it to feedback
composition.

At the interface level, the plant is a Mealy process
\[
P:A\rightsquigarrow O\times R.
\]
The adaptive controller is a Mealy process
\[
K:O\times R\rightsquigarrow A\times V.
\]
The action output of \(K\) is connected to the action input of \(P\), and the observation/reward
output of \(P\) is connected to the observation/reward input of \(K\).

Equivalently, in signed-interface notation, the plant may be viewed as a component
\[
P:\mathbf 1\longrightarrow (O\times R,A),
\]
because it produces observation/reward data and consumes action data.  The adaptive controller may
be viewed as a component
\[
K:(O\times R,A)\longrightarrow (V,0),
\]
because it consumes observation/reward data, produces action data, and exposes a value output.

Their composite is a closed-loop component
\[
\mathbf 1\longrightarrow (V,0).
\]
The interface
\[
(O\times R,A)
\]
is hidden by feedback.

Thus the trace or Int construction supplies the feedback wiring:
\[
\boxed{
\text{trace/Int supplies the signed feedback discipline}.
}
\]
The causal closed-loop Mealy transition supplies the state-mediated adaptive dynamics:
\[
\boxed{
\text{causal Mealy execution supplies the learning dynamics}.
}
\]

The optimizer and updater are not created by the trace.  They are stateful components of the
adaptive controller.  The trace closes the external plant-controller loop.

\section{Specifications of Deterministic Cybernetic Systems}
\label{sec:cybernetic-specifications}

Once the closed-loop system is represented by a transition
\[
F:S\longrightarrow S,
\]
a state specification is a subobject
\[
\varphi\hookrightarrow S.
\]
In \(\mathbf{Set}\), this is simply a subset \(\varphi\subseteq S\).

The one-step diamond modality is
\[
\langle F\rangle\varphi.
\]
Since \(F\) is deterministic and functional, this is simply the inverse image
\[
F^{-1}(\varphi).
\]
Likewise, the box modality
\[
[F]\varphi
\]
also reduces to inverse image in the deterministic functional case.

For iterated execution, we write
\[
\langle F^*\rangle \varphi
\]
for eventual reachability of \(\varphi\), and
\[
[F^*]\varphi
\]
for invariance of \(\varphi\) along all finite executions.

Typical cybernetic specifications include:
\[
\langle F^*\rangle\mathrm{Comfort}
\]
for eventual comfort in a thermostat,
\[
[F^*]\mathrm{SafeTemp}
\]
for temperature safety,
\[
\langle F^*\rangle(\theta_b>\theta_a)
\]
for eventual learning of a better bandit arm,
\[
[F^*]\mathrm{NoStockout}
\]
for inventory safety,
\[
\langle F^*\rangle\mathrm{Target}
\]
for navigation reachability, and
\[
[F^*](d>0)
\]
for collision avoidance in cruise control.

In span-valued or nondeterministic variants, diamonds and boxes diverge into their usual may/must
interpretations.  In the deterministic examples below, they coincide at one step but remain
conceptually useful for reachability and invariance.

\section{Simulations and Abstraction of Closed-Loop Systems}
\label{sec:cybernetic-simulations}

Closed-loop cybernetic systems can be compared by simulations.

\begin{definition}[Strict closed-loop simulation]
\label{def:strict-closed-loop-simulation}
Let
\[
F:S\longrightarrow S
\]
and
\[
F':S'\longrightarrow S'
\]
be deterministic closed-loop systems.  A strict simulation from \(F\) to \(F'\) is a map
\[
\sigma:S\longrightarrow S'
\]
such that
\[
\sigma(F(s))=F'(\sigma(s))
\]
for every \(s\in S\).
Equivalently, the square
\[
\begin{tikzcd}
S \arrow[r,"F"] \arrow[d,"\sigma"'] & S \arrow[d,"\sigma"] \\
S' \arrow[r,"F'"'] & S'
\end{tikzcd}
\]
commutes.
\end{definition}

A simulation may compare plants, controllers, optimizers, updaters, or entire closed-loop systems.
For example:
\[
\text{a plant simulation abstracts one plant by another},
\]
\[
\text{a controller simulation compares two controllers},
\]
\[
\text{an optimizer simulation preserves selected actions and values},
\]
\[
\text{a closed-loop simulation compares whole cybernetic systems}.
\]

\begin{proposition}[Preservation of forward behaviour]
\label{prop:simulation-forward-behaviour}
If \(\sigma:S\to S'\) is a strict simulation from \(F\) to \(F'\), then for every \(n\geq 0\),
\[
\sigma(F^n(s))=(F')^n(\sigma(s)).
\]
\end{proposition}

\begin{proof}
The proof is by induction on \(n\).  The case \(n=0\) is immediate.  If the statement holds for
\(n\), then
\[
\sigma(F^{n+1}(s))
=
\sigma(F(F^n(s)))
=
F'(\sigma(F^n(s)))
=
F'((F')^n(\sigma(s)))
=
(F')^{n+1}(\sigma(s)).
\]
\end{proof}

Thus simulations preserve forward trajectories.  Consequently, specifications expressed at the
abstract level may be transported along suitable simulation maps.  In span form, one replaces the
commuting square by a compatible simulation span and obtains the corresponding relational version.

\section{Example I: Adaptive Thermostat}
\label{sec:example-adaptive-thermostat}

The first example is a finite adaptive thermostat.

Let the plant state be the temperature level
\[
X=\{0,1,2,3,4\}.
\]
Let the actions be
\[
A=\{H,I,C\},
\]
where \(H\) means heat, \(I\) means idle, and \(C\) means cool.  The plant transition is
\[
p(x,H)=\min(4,x+1),
\]
\[
p(x,I)=x,
\]
\[
p(x,C)=\max(0,x-1).
\]
The observation is full-state observation:
\[
h(x)=x.
\]

Let the controller state record the previous action:
\[
C=A.
\]
Let the policy state be a learned comfort target:
\[
\Theta=\{2,3\}.
\]
Let the optimizer state be a tie-breaking or cache state:
\[
Z=A.
\]
Let the updater memory be the previous observation:
\[
U=X.
\]

For \(\epsilon>0\) small, define the evaluator
\[
Q(\theta,c,o,a)
=
-\left|p(o,a)-\theta\right|
-\epsilon\cdot \mathbf 1_{a\neq c}.
\]
The first term rewards moving toward the target.  The second term penalizes switching.

The optimizer returns
\[
(a^*,v^*,z^+)=\Omega(z,\theta,c,o),
\]
where
\[
a^*\in \operatorname{argmax}_{a\in A}Q(\theta,c,o,a),
\]
\[
v^*=\max_{a\in A}Q(\theta,c,o,a),
\]
and
\[
z^+=a^*.
\]
If there are exact ties, fix the deterministic priority
\[
I\succ H\succ C.
\]

The updater is
\[
\theta'=
\begin{cases}
2 & \text{if } o'\leq 1,\\
3 & \text{if } o'\geq 4,\\
\theta & \text{otherwise},
\end{cases}
\]
and
\[
c'=a^*,
\qquad
z'=z^+,
\qquad
u'=o'.
\]

The closed-loop state is
\[
S=X\times A\times \{2,3\}\times A\times X.
\]

\subsection{Trajectory from a cold state}
\label{subsec:thermostat-cold-trajectory}

Take initial state
\[
(x,c,\theta,z,u)=(0,I,3,I,0).
\]
Then one possible trajectory is:

\[
\begin{array}{c|c|c|c|c|c|c}
t & x_t & c_t & \theta_t & a_t^* & x_{t+1} & v_t^* \\
\hline
0 & 0 & I & 3 & H & 1 & -2-\epsilon \\
1 & 1 & H & 2 & H & 2 & 0 \\
2 & 2 & H & 2 & I & 2 & -\epsilon \\
3 & 2 & I & 2 & I & 2 & 0 \\
4 & 2 & I & 2 & I & 2 & 0
\end{array}
\]

The system moves from too cold to the comfort region and then remains there.

\subsection{Trajectory from a hot state}
\label{subsec:thermostat-hot-trajectory}

Take initial state
\[
(x,c,\theta,z,u)=(4,I,2,I,4).
\]
Then:

\[
\begin{array}{c|c|c|c|c|c|c}
t & x_t & c_t & \theta_t & a_t^* & x_{t+1} & v_t^* \\
\hline
0 & 4 & I & 2 & C & 3 & -1-\epsilon \\
1 & 3 & C & 3 & I & 3 & -\epsilon \\
2 & 3 & I & 3 & I & 3 & 0 \\
3 & 3 & I & 3 & I & 3 & 0
\end{array}
\]

The example exhibits all components:
\[
X=\text{plant temperature},
\]
\[
C=\text{previous action},
\]
\[
\Theta=\text{learned comfort target},
\]
\[
Z=\text{optimizer tie-breaking/cache state},
\]
\[
U=\text{previous observation}.
\]

The argmax output controls the plant.  The max output records the quality of the best currently
available control action.

\section{Example II: Deterministic Bandit Learning}
\label{sec:example-deterministic-bandit}

The deterministic bandit is the smallest example in which learning is visible.

Let
\[
A=\{a,b\}
\]
be the action set.  The plant state is a deterministic reward table
\[
X=\{0,1\}^A.
\]
Fix the concrete plant
\[
x(a)=0,
\qquad
x(b)=1.
\]
The plant state does not change:
\[
p(x,i)=x.
\]
The reward is
\[
\rho(x,i,x)=x(i).
\]
The observation object is terminal:
\[
O=1.
\]

The policy state is an estimated reward table
\[
\Theta=\{0,1\}^A.
\]
Write
\[
\theta=(\hat r_a,\hat r_b).
\]
Let the optimizer state be
\[
Z=\{\operatorname{try}\,a,\operatorname{try}\,b,\operatorname{exploit}\}.
\]
Let the updater memory be
\[
U=A+\{\bot\}.
\]

The evaluator is
\[
Q(\theta,i)=\hat r_i.
\]
The optimizer is:
\[
\Omega(\operatorname{try}\,a,\theta)
=
(a,\hat r_a,\operatorname{try}\,b),
\]
\[
\Omega(\operatorname{try}\,b,\theta)
=
(b,\hat r_b,\operatorname{exploit}),
\]
and
\[
\Omega(\operatorname{exploit},\theta)
=
(i^*,v^*,\operatorname{exploit}),
\]
where
\[
i^*\in \operatorname{argmax}_{i\in A}\hat r_i,
\qquad
v^*=\max_{i\in A}\hat r_i.
\]
Ties are broken by the priority \(a\succ b\).

The updater sets the estimate of the selected action to the observed reward:
\[
\hat r_i'=r
\]
for the chosen action \(i\), and leaves the other estimate fixed.

Starting with
\[
\theta_0=(0,0),
\qquad
z_0=\operatorname{try}\,a,
\]
the trajectory is:

\[
\begin{array}{c|c|c|c|c|c|c}
t & z_t & \theta_t & i_t & r_t & \theta_{t+1} & v_t^* \\
\hline
0 & \operatorname{try}\,a & (0,0) & a & 0 & (0,0) & 0 \\
1 & \operatorname{try}\,b & (0,0) & b & 1 & (0,1) & 0 \\
2 & \operatorname{exploit} & (0,1) & b & 1 & (0,1) & 1 \\
3 & \operatorname{exploit} & (0,1) & b & 1 & (0,1) & 1
\end{array}
\]

This example separates exploration, optimization, and update.  The optimizer chooses using the
current estimates.  The updater changes the estimates after observing reward.  Therefore the
optimizer is not the learner; the updater is the learner.

\section{Example III: Finite Inventory Regulation}
\label{sec:example-inventory-regulation}

We next consider a finite deterministic inventory controller.

Let inventory levels be
\[
I=\{0,1,2,3,4,5\}.
\]
Let demand phase be
\[
D=\{0,1,2\}.
\]
The plant state is
\[
X=I\times D.
\]
Demand is deterministic and cyclic:
\[
d_0=1,
\qquad
d_1=3,
\qquad
d_2=2.
\]
Actions are order quantities:
\[
A=\{0,1,2,3\}.
\]

If the plant state is \((x,i)\) and the order is \(a\), define
\[
y=\min(5,x+a).
\]
The resulting inventory is
\[
x'=\max(0,y-d_i),
\]
and the demand phase advances by
\[
i'=i+1\pmod 3.
\]
Thus
\[
p((x,i),a)=(x',i').
\]
The observation is
\[
h(x,i)=(x,d_i).
\]

A simple performance signal is negative cost.  Let shortage be
\[
\operatorname{short}(x,i,a)=\max(0,d_i-\min(5,x+a)),
\]
and holding be
\[
\operatorname{hold}(x,i,a)=\max(0,\min(5,x+a)-d_i).
\]
Define
\[
\rho((x,i),a,(x',i'))
=
-\bigl(2\operatorname{short}(x,i,a)+\operatorname{hold}(x,i,a)+a\bigr).
\]

Let the controller state be a cooldown bit:
\[
C=\{0,1\}.
\]
Let the policy state be a demand estimate:
\[
\Theta=\{1,2,3\}.
\]
Let the optimizer state store the previous selected order:
\[
Z=A.
\]
Let updater memory store the previous observed demand:
\[
U=\{1,2,3\}.
\]

The guard is
\[
G(c,o)=
\begin{cases}
\{0,1\} & \text{if } c=1,\\
\{0,1,2,3\} & \text{if } c=0.
\end{cases}
\]
Thus after a large order, the controller temporarily restricts the next order.

For observation \(o=(x,d_i)\), define the evaluator
\[
Q(\theta,c,o,a)
=
-\left|x+a-\theta\right|-a.
\]
The optimizer chooses
\[
a^*\in \operatorname{argmax}_{a\in G(c,o)}Q(\theta,c,o,a),
\]
and returns
\[
v^*=\max_{a\in G(c,o)}Q(\theta,c,o,a).
\]

The updater sets
\[
\theta'=d_i,
\]
and
\[
c'=
\begin{cases}
1 & \text{if } a^*=3,\\
0 & \text{otherwise}.
\end{cases}
\]
It also sets
\[
z'=a^*,
\qquad
u'=d_i.
\]

A sample trajectory beginning at
\[
((x,i),c,\theta,z,u)=((1,0),0,1,0,1)
\]
is:

\[
\begin{array}{c|c|c|c|c|c|c|c|c}
t & x_t & i_t & d_{i_t} & c_t & \theta_t & a_t^* & v_t^* & x_{t+1} \\
\hline
0 & 1 & 0 & 1 & 0 & 1 & 0 & 0 & 0 \\
1 & 0 & 1 & 3 & 0 & 1 & 1 & -1 & 0 \\
2 & 0 & 2 & 2 & 0 & 3 & 2 & -3 & 0 \\
3 & 0 & 0 & 1 & 0 & 2 & 1 & -1 & 0 \\
4 & 0 & 1 & 3 & 0 & 1 & 1 & -1 & 0
\end{array}
\]

The controller is not optimal in any sophisticated operations-research sense.  That is not the point.
The point is structural: the plant state, controller guard state, demand-estimate state, optimizer
state, and updater memory have distinct roles.

\section{Example IV: Finite Navigation with Learned Map}
\label{sec:example-navigation-learned-map}

This example illustrates model learning.

Let positions be
\[
P=\{0,1,2,3\},
\]
with target position \(3\).  Let actions be
\[
A=\{L,S,R\},
\]
where \(L\) moves left, \(S\) stays, and \(R\) moves right.

The plant state includes a wall bit:
\[
X=P\times\{0,1\}.
\]
The value \(w=1\) means that the edge \(1\to 2\) is blocked.  The wall bit is part of the plant
state, but the controller may not initially know it.

For \(w\in\{0,1\}\), define:
\[
p((q,w),L)=(\max(0,q-1),w),
\]
\[
p((q,w),S)=(q,w),
\]
and
\[
p((q,w),R)=
\begin{cases}
(1,w) & \text{if } q=1 \text{ and } w=1,\\
(\min(3,q+1),w) & \text{otherwise}.
\end{cases}
\]
The observation is
\[
o=(q,\beta),
\]
where \(\beta\in\{0,1\}\) records whether the last action produced a bump.  In this finite
presentation, the bump signal may be included in updater memory, or equivalently in the post-action
observation.

Let the policy/model state be
\[
\Theta=\{0,1,?\}.
\]
The value \(?\) means that the controller has not yet learned whether the wall is present.

Let the controller mode be
\[
C=\{\operatorname{explore},\operatorname{avoid},\operatorname{done}\}.
\]
Let the optimizer state and updater memory be finite sets storing the previous action and bump
signal:
\[
Z=A,
\qquad
U=A\times\{0,1\}.
\]

The evaluator uses the learned model to predict progress toward the target:
\[
Q(\theta,c,o,a)
=
-\operatorname{dist}_{\theta}(p_{\theta}(o,a),3)
-
\operatorname{risk}_{\theta}(o,a).
\]
Here \(p_\theta\) is the controller's learned transition model, and \(\operatorname{risk}_\theta\)
penalizes attempting \(R\) from position \(1\) when the wall status is unknown or known blocked.

The optimizer chooses
\[
a^*\in\operatorname{argmax}_{a\in A}Q(\theta,c,o,a),
\]
and returns
\[
v^*=\max_{a\in A}Q(\theta,c,o,a).
\]

The updater modifies the learned wall state:
\[
\theta'=
\begin{cases}
1 & \text{if action } R \text{ from position } 1 \text{ produces a bump},\\
0 & \text{if action } R \text{ from position } 1 \text{ succeeds},\\
\theta & \text{otherwise}.
\end{cases}
\]
The controller mode updates by
\[
c'=
\begin{cases}
\operatorname{done} & \text{if } q'=3,\\
\operatorname{avoid} & \text{if } \theta'=1,\\
\operatorname{explore} & \text{otherwise}.
\end{cases}
\]

If the wall is absent, a typical trajectory is:
\[
\begin{array}{c|c|c|c|c|c}
t & q_t & \theta_t & c_t & a_t^* & q_{t+1} \\
\hline
0 & 0 & ? & \operatorname{explore} & R & 1 \\
1 & 1 & ? & \operatorname{explore} & R & 2 \\
2 & 2 & 0 & \operatorname{explore} & R & 3 \\
3 & 3 & 0 & \operatorname{done} & S & 3
\end{array}
\]

If the wall is present, the trajectory begins:
\[
\begin{array}{c|c|c|c|c|c|c}
t & q_t & \theta_t & c_t & a_t^* & \text{bump} & \theta_{t+1} \\
\hline
0 & 0 & ? & \operatorname{explore} & R & 0 & ? \\
1 & 1 & ? & \operatorname{explore} & R & 1 & 1 \\
2 & 1 & 1 & \operatorname{avoid} & S & 0 & 1
\end{array}
\]

The example shows that deterministic learning need not be reward learning.  It can be model
learning.  The policy/model state \(\Theta\) stores an internal approximation to the plant.

The natural specifications are:
\[
\langle F^*\rangle\operatorname{Target}
\]
for eventual target reachability, and
\[
[F^*]\operatorname{Safe}
\]
for safety.  In the wall-present case, reachability may fail unless additional actions or paths are
available.  The formalism exposes this failure as a property of the closed-loop transition.

\section{Example V: Adaptive Cruise Control}
\label{sec:example-adaptive-cruise}

We now consider a finite deterministic cruise-control system.

Let following distance be
\[
d\in\{0,1,2,3,4\}.
\]
Let ego velocity be
\[
v\in\{0,1,2\}.
\]
Let lead-car velocity be
\[
\ell\in\{0,1,2\}.
\]
The plant state is
\[
X=\{0,1,2,3,4\}\times\{0,1,2\}\times\{0,1,2\}.
\]
Actions are accelerations:
\[
A=\{-1,0,+1\}.
\]

Given state \((d,v,\ell)\) and action \(a\), the ego velocity updates by
\[
v'=\min(2,\max(0,v+a)).
\]
The following distance updates by
\[
d'=\min(4,\max(0,d+\ell-v')).
\]
For simplicity, keep the lead-car velocity constant:
\[
\ell'=\ell.
\]
Thus
\[
p((d,v,\ell),a)=(d',v',\ell').
\]
The observation is full:
\[
h(d,v,\ell)=(d,v,\ell).
\]

The reward is a negative safety/performance cost:
\[
\rho((d,v,\ell),a,(d',v',\ell'))
=
-\left(
10\cdot\mathbf 1_{d'=0}
+
|v'-\ell|
+
\mathbf 1_{a=-1}
\right).
\]

Let the policy state be the learned safe following gap:
\[
\Theta=\{1,2,3\}.
\]
Let the controller mode be
\[
C=\{\operatorname{normal},\operatorname{braking}\}.
\]
Let \(Z=A\) store the previous optimizer choice, and let \(U\) store the previous observation.

The evaluator is
\[
Q(\theta,c,(d,v,\ell),a)
=
-|v'-\ell|
-
10\cdot\mathbf 1_{d'<\theta}
-
\mathbf 1_{a=-1},
\]
where \(v'\) and \(d'\) are computed from the plant prediction under action \(a\).

The optimizer returns
\[
a^*\in\operatorname{argmax}_{a\in A}Q(\theta,c,o,a),
\]
and
\[
v^*=\max_{a\in A}Q(\theta,c,o,a).
\]

The updater changes the safe-gap estimate by
\[
\theta'=
\begin{cases}
\min(3,\theta+1) & \text{if } d'\leq \theta,\\
\max(1,\theta-1) & \text{if } d'\geq \theta+2,\\
\theta & \text{otherwise},
\end{cases}
\]
and changes the controller mode by
\[
c'=
\begin{cases}
\operatorname{braking} & \text{if } d'\leq \theta,\\
\operatorname{normal} & \text{if } d'>\theta.
\end{cases}
\]
Set
\[
z'=a^*,
\qquad
u'=o'.
\]

A safe initial trajectory might begin:
\[
\begin{array}{c|c|c|c|c|c|c|c}
t & d_t & v_t & \ell_t & \theta_t & c_t & a_t^* & v_t^* \\
\hline
0 & 4 & 1 & 1 & 2 & \operatorname{normal} & 0 & 0 \\
1 & 4 & 1 & 1 & 1 & \operatorname{normal} & 0 & 0 \\
2 & 4 & 1 & 1 & 1 & \operatorname{normal} & 0 & 0
\end{array}
\]

An unsafe initial trajectory might begin:
\[
\begin{array}{c|c|c|c|c|c|c|c}
t & d_t & v_t & \ell_t & \theta_t & c_t & a_t^* & v_t^* \\
\hline
0 & 1 & 2 & 0 & 2 & \operatorname{normal} & -1 & -11 \\
1 & 0 & 1 & 0 & 3 & \operatorname{braking} & -1 & -11 \\
2 & 0 & 0 & 0 & 3 & \operatorname{braking} & 0 & -10
\end{array}
\]

The precise numerical values are less important than the structural separation:
\[
\text{plant dynamics},
\quad
\text{safety mode},
\quad
\text{adaptive safe-gap parameter},
\quad
\text{optimizer},
\quad
\text{updater}.
\]
The specification
\[
[F^*](d>0)
\]
expresses collision avoidance.  The example shows that an optimizer can choose the best available
action while still exposing a low value certificate \(v^*\), indicating that every admissible action
is bad from the current state.

\section{Example VI: Two Concurrent Plants Under One Adaptive Controller}
\label{sec:example-concurrent-plants}

This example explains why the causal update discipline does not collapse concurrency into
interleaving.

Let two independent rooms be controlled by a shared energy manager.  For each \(i=1,2\), let
\[
X_i=\{0,1,2,3,4\}
\]
be the temperature state of room \(i\), and let
\[
A_i=\{H,I,C\}
\]
be its action set.  The individual plant dynamics are
\[
p_i(x_i,H)=\min(4,x_i+1),
\]
\[
p_i(x_i,I)=x_i,
\]
\[
p_i(x_i,C)=\max(0,x_i-1).
\]
The two plants combine in parallel:
\[
X=X_1\times X_2,
\qquad
A=A_1\times A_2,
\]
with transition
\[
p((x_1,x_2),(a_1,a_2))
=
(p_1(x_1,a_1),p_2(x_2,a_2)).
\]
This product transition is not an interleaving.  The two plant coordinates update independently.

The controller observes
\[
o=(x_1,x_2).
\]
Let the policy state be target temperatures
\[
\Theta=\{2,3\}\times\{2,3\}.
\]
Let the controller state \(C\) store an energy mode:
\[
C=\{\operatorname{normal},\operatorname{saving}\}.
\]

The admissible joint actions are constrained by a shared energy budget:
\[
G(c,o)
=
\{(a_1,a_2)\in A_1\times A_2\mid \text{not both }a_1=H\text{ and }a_2=H\}.
\]
In saving mode, one may impose the stronger guard
\[
G(\operatorname{saving},o)
=
\{(a_1,a_2)\mid a_1\neq H,\ a_2\neq H\}.
\]

The evaluator is
\[
Q(\theta,c,o,(a_1,a_2))
=
-\left|p_1(x_1,a_1)-\theta_1\right|
-\left|p_2(x_2,a_2)-\theta_2\right|
-\operatorname{energy}(a_1,a_2).
\]
The optimizer chooses the best admissible joint action:
\[
(a_1^*,a_2^*)\in
\operatorname{argmax}_{(a_1,a_2)\in G(c,o)}
Q(\theta,c,o,(a_1,a_2)),
\]
and returns the corresponding maximum value.

The important point is that the plant side is genuinely parallel:
\[
P_1\times P_2:
A_1\times A_2
\rightsquigarrow
O_1\times O_2.
\]
The controller may introduce a coupled optimization problem, but this does not mean the plant
updates are interleaved.  Causal dependence is not interleaving semantics.

Thus this deterministic cybernetic pattern remains compatible with product-parallel and cubical
true-concurrency structure.  The update depends on dataflow, not on a global sequential scheduler.

\section{Example VII: Optimizer Comparison and Value Certificates}
\label{sec:example-optimizer-comparison}

The optimizer was required to output both
\[
\operatorname{argmax}
\]
and
\[
\max.
\]
This example explains why.

Consider a plant and evaluator for which the selected action controls the plant, but the value of
the selected action is also observed by a monitor.  The action output is
\[
a^*,
\]
whereas the value certificate is
\[
v^*.
\]

For instance, in the cruise-control example,
\[
a^*
\]
is the acceleration command, while
\[
v^*
\]
measures the value of the best available acceleration.  A low \(v^*\) says that the controller is
in a bad situation even after optimizing.

This supports specifications such as
\[
[F^*](v^*\geq \alpha),
\]
which says that the closed-loop system always has an available control action of value at least
\(\alpha\).  It also supports
\[
\langle F^*\rangle(v^*<\alpha),
\]
which says that degraded control is eventually reachable.

One may also express fallback policies:
\[
v^*<\alpha
\Longrightarrow
\operatorname{EmergencyMode}.
\]

Compare three controllers:
\[
K_{\operatorname{act}},
\qquad
K_{\operatorname{val}},
\qquad
K_{\operatorname{both}}.
\]
The first exposes only the selected action \(a^*\).  The second exposes only the value \(v^*\).  The
third exposes both:
\[
(a^*,v^*).
\]
Only the third simultaneously controls the plant and certifies the quality of the current decision.

Therefore the optimizer should not be represented only by an action map
\[
O\longrightarrow A.
\]
It should be represented by a decision-value map
\[
O\longrightarrow A\times W,
\]
or, in the stateful setting,
\[
\Omega:Z\times\Theta\times C\times O\longrightarrow A\times W\times Z.
\]

\section{Behaviour, Eventual Periodicity, and Black-Boxing}
\label{sec:cybernetic-behaviour-black-boxing}

Let the closed-loop transition be
\[
F:S\longrightarrow S.
\]
Let the visible output of the closed system be
\[
b:S\longrightarrow B.
\]
For example,
\[
B=O\times A\times R\times W
\]
may record observation, selected action, reward, and optimizer value.

The visible behaviour from an initial state \(s_0\in S\) is the stream
\[
b(s_0),\ b(F(s_0)),\ b(F^2(s_0)),\ldots.
\]

\begin{proposition}[Eventually periodic visible behaviour]
\label{prop:eventually-periodic-visible-behaviour}
If \(S\) is finite, then for every initial state \(s_0\), the visible behaviour stream
\[
b(s_0),\ b(F(s_0)),\ b(F^2(s_0)),\ldots
\]
is eventually periodic.
\end{proposition}

\begin{proof}
Since \(S\) is finite, the state sequence
\[
s_0,\ F(s_0),\ F^2(s_0),\ldots
\]
eventually repeats.  Applying \(b\) to this eventually periodic sequence gives an eventually
periodic output sequence.
\end{proof}

This gives a useful black-box perspective on the examples.  The thermostat stabilizes or cycles.
The deterministic bandit eventually exploits.  The inventory controller may enter a demand-cycle
orbit.  The navigation system reaches the target or loops.  The cruise controller stabilizes,
oscillates, or enters emergency behaviour.

Thus the internal plant/controller/policy/optimizer/updater states may be hidden, leaving only the
observable closed-loop trace.

\section{What the Examples Establish}
\label{sec:what-examples-establish}

The examples justify the formal pattern.  In each case, the closed-loop state has the form
\[
S=X\times C\times\Theta\times Z\times U.
\]
The plant state \(X\) belongs to the regulated system.  The controller state \(C\) records modes,
guards, or safety information.  The policy/model state \(\Theta\) records learned predictive or
evaluative data.  The optimizer state \(Z\) records tie-breaking, search, or cache data.  The
updater state \(U\) records the memory needed to change the future policy/model state.

Each example also uses:
\[
Q:\Theta\times C\times O\times A\longrightarrow W
\]
as evaluator,
\[
\Omega:Z\times\Theta\times C\times O\longrightarrow A\times W\times Z
\]
as optimizer, and
\[
L:
C\times\Theta\times Z\times U
\times O\times A\times W\times R\times O
\longrightarrow
C\times\Theta\times Z\times U
\]
as updater.

The main distinctions are:

\[
\begin{array}{c|l}
\text{distinction} & \text{meaning} \\
\hline
X \text{ versus } C & \text{plant dynamics versus controller mode} \\
\Theta \text{ versus } \pi_\theta & \text{policy state versus policy readout} \\
Q \text{ versus } \Omega & \text{evaluation versus selection} \\
\Omega \text{ versus } L & \text{current decision versus future learning} \\
a^* \text{ versus } v^* & \text{control action versus value certificate} \\
\text{trace/Int} \text{ versus } F & \text{feedback wiring versus closed-loop execution}
\end{array}
\]

The flat systems diagram is therefore repaired as follows:
\[
\boxed{
\text{plant and adaptive controller are feedback-composed through signed interfaces;}
}
\]
\[
\boxed{
\text{the adaptive controller internalizes policy state, optimization, and update;}
}
\]
\[
\boxed{
\text{the resulting closed-loop system is deterministic Markov on the coupled state object.}
}
\]

\section{Chapter Summary}
\label{sec:cybernetic-summary}

This chapter showed that deterministic cybernetics is a natural application of categorical automata
theory.

First, ordinary deterministic Mealy machines were identified with Dirac-valued controlled Markov
processes.  A Mealy transition-output pair
\[
\tau:S\times I\to S,
\qquad
\omega:S\times I\to O
\]
determines the Dirac kernel
\[
s\longmapsto \delta_{(\tau(s,i),\omega(s,i))}.
\]

Second, plants were represented as stateful Mealy systems
\[
P:A\rightsquigarrow O\times R.
\]
They receive actions, update plant state, and emit observations and performance signals.

Third, adaptive controllers were represented as stateful Mealy systems
\[
K:O\times R\rightsquigarrow A\times V
\]
whose internal state decomposes as
\[
C\times\Theta\times Z\times U.
\]
This decomposition separates controller mode, policy/model state, optimizer state, and updater
memory.

Fourth, the optimizer was represented as a stateful argmax/max component
\[
\Omega:Z\times\Theta\times C\times O\to A\times W\times Z.
\]
Its action projection controls the plant.  Its value projection certifies the quality of the current
decision.

Fifth, the updater was represented as the learning component
\[
L:
C\times\Theta\times Z\times U
\times O\times A\times W\times R\times O
\to
C\times\Theta\times Z\times U.
\]
It changes the future decision problem after observing the plant consequence.

Sixth, the plant and adaptive controller were closed into a single deterministic transition
\[
F:X\times C\times\Theta\times Z\times U
\to
X\times C\times\Theta\times Z\times U.
\]
The associated Markov process is the Dirac kernel
\[
s\longmapsto \delta_{F(s)}.
\]

Finally, the examples showed that the same pattern applies to thermostatic regulation, deterministic
bandit learning, inventory control, navigation with learned maps, adaptive cruise control, concurrent
plants under a shared controller, and optimizer-value monitoring.

Thus deterministic cybernetics is not an external metaphor added to categorical automata theory.
It is the closed-loop adaptive-control fragment of the theory: stateful open systems connected by
feedback, equipped with optimization, update, behaviour, and specification.